\documentclass[10pt,a4paper]{amsart}
\usepackage[margin=19mm]{geometry}
\usepackage{amsmath,amssymb,mathtools}
\usepackage[scr=rsfs,cal=euler]{mathalfa}
\usepackage{xfrac}
\usepackage[unicode,hidelinks,bookmarks=false]{hyperref}
\newcommand{\ie}{i.e.\ }
\newcommand{\cf}{cf.\ }
\newcommand{\resp}{respectively\ }
\newcommand{\Subsection}{\subsection}
\providecommand{\nobreakdash}{\nobreak\hbox{-}}
\setlength{\emergencystretch}{2em}
\multlinegap0pt
\usepackage{amssymb}
\usepackage{mathtools}
\newtheorem{proposition}{Proposition}
\title{An incomplete attack on the upper bound of the unit distance problem}
\author{Steven Senger}
\date{Source: arXiv:2605.26145v1 (CC BY 4.0)}
\begin{document}
\maketitle

\noindent\emph{Source and licence.} An incomplete attack on the upper bound of the unit distance problem. Version arXiv:2605.26145v1 (CC BY 4.0), distributed by arXiv under the Creative Commons Attribution 4.0 International licence, http://creativecommons.org/licenses/by/4.0/, as recorded in the arXiv licence metadata for this exact version and shown on the abstract page https://arxiv.org/abs/2605.26145v1. Full licence notice: \texttt{licenses/arxiv-2605.26145-CC-BY-4.0.txt}.\par\smallskip

\noindent\textit{Editorial scope.} Counted unit: the article's proposition distArc (source0.tex lines 142--148). The article's complete original proof of it is delivered with it (source0.tex lines 149--169, one \texttt{proof} environment of 21 lines). No mathematics is written, repaired or re-derived: the statement and the proof are the article's own lines, in the article's own notation, and no retained line is substituted or re-flowed.  No further source-local context is needed: every cross-reference of every retained line resolves inside the two delivered spans.  Nothing was omitted from any retained span, so the package carries no omission record.  No retained line contains a citation, so the wrapper carries no bibliography and no bibliography entry.  No retained line contains a figure environment, an image-inclusion command or drawing code, so no image is delivered and the package declares no asset dependency.  Editorial formatting only, in addition: an AMS \texttt{amsart} preamble that restates the article's own declarations and macros used by the retained lines, namely the environments \texttt{proposition} on separate counters, together with the standard packages those lines need.  The retained environments print in this document as Proposition 1 in order of appearance, whereas the article prints the same items with the numbering of its own full document.  The complete native source of the article as distributed in the arXiv e-print for this exact version is bundled unchanged as \texttt{sources/201399/source0.tex}.
\begin{proposition}\label{distArc}
Given distinct points $a$ and $b$ on a unit circle separated by an arc of length $\theta <\frac{\pi}{2}:$
\begin{enumerate}
\item[(i)] The distance from $a$ to $b$ is $\sqrt{2-2\cos\theta}.$
\item[(ii)] $\theta - \frac{\theta^2}{\sqrt{12}} < \sqrt{2-2\cos\theta} < \theta - \frac{\theta^3}{25}$
\end{enumerate}
\end{proposition}

\begin{proof}
$(i)$ The arc of the circle subtends the angle $\theta,$ so we can directly apply the law of cosines to the triangle formed by $a, b,$ and the center of the circle.

$(ii)$
With this in mind, estimate from below by
$$
\sqrt{2-2\cos \theta} \geq \sqrt{2-2\left(1-\frac{\theta^2}{2!}+\frac{\theta^4}{4!}\right)}=\sqrt{\theta^2-\frac{\theta^4}{12}}
$$
$$
=\sqrt{\left(\theta-\frac{\theta^2}{\sqrt{12}}\right)\left(\theta+\frac{\theta^2}{\sqrt{12}}\right)}> \sqrt{\left(\theta-\frac{\theta^2}{\sqrt{12}}\right)^2}.
$$
Similarly, we estimate from above by
$$
\sqrt{2-2\cos \theta} \leq \sqrt{2-2\left(1-\frac{\theta^2}{2!}+\frac{\theta^4}{4!}-\frac{\theta^6}{6!}\right)}
=\sqrt{\theta^2-\frac{\theta^4}{12}+\frac{\theta^6}{360}}
$$
$$
< \sqrt{\theta^2-\frac{2\cdot\theta^4}{25}+\frac{\theta^6}{625}} = \sqrt{\left(\theta - \frac{\theta^3}{25}\right)^2},
$$
where in the last line, we used the fact that $\theta<\pi/2< \sqrt{2250/795}.$
\end{proof}

\end{document}
